%! Inverses of Exponential Functions — Unit 2, Lesson 2.10 — Homework (Answer Key)
\documentclass[10pt]{article}
\usepackage{precalculus-article}
\usepackage{precalculus-key}   % pulls in -boxes; do NOT also load -boxes
\newcommand{\TallMath}[1]{$\displaystyle #1\rule[-1.4em]{0pt}{3.2em}$}

\begin{document}

\pageheader{Unit 2, Lesson 2.10}{Homework --- Answer Key}
\namedateperiod

\vspace{0.1in}

\begin{enumerate}[leftmargin=1.8em, itemsep=14pt]

\item Let $g(x) = 6^x$.

\begin{enumerate}[leftmargin=1.5em, itemsep=6pt, label=\alph*.]
  \item Write the inverse function $f(x)$.

  \vspace{4pt}
  $f(x) = $ \ans{$\log_6 x$}

  \item Complete the table and inverse pairs.

  \vspace{4pt}
  \begin{center}
  \renewcommand{\arraystretch}{1.8}
  \begin{tabular}{|c|c|c|}
  \hline
  \rowcolor{sky}
  $x$ & $g(x) = 6^x$ & Point on $f$: $(g(x),\; x)$ \\
  \hline
  $0$ & \ans{1}  & \ans{$(1,\;0)$} \\
  \hline
  $1$ & \ans{6}  & \ans{$(6,\;1)$} \\
  \hline
  $2$ & \ans{36} & \ans{$(36,\;2)$} \\
  \hline
  \end{tabular}
  \end{center}
\end{enumerate}

\item Verify $f(x) = \log_5 x$ and $g(x) = 5^x$ are inverses.

\vspace{4pt}
$f(g(x)) = \log_5(5^x) = $ \ans{$x$}

\vspace{6pt}
$g(f(x)) = 5^{\log_5 x} = $ \ans{$x$}

\vspace{4pt}
Conclusion: \ansline{Since both compositions equal $x$, the functions $f$ and $g$ are inverses
of each other.}

\item Algae model $A(t) = 10^t$.

\begin{enumerate}[leftmargin=1.5em, itemsep=6pt, label=\alph*.]
  \item $t(A) = $ \ans{$\log_{10} A$}

  \item $t(1000) = \log_{10} 1000 = $ \ans{3}

  \vspace{4pt}
  In context: \ansline{After 3 days, the algae covers 1000 cm$^2$; equivalently, it takes 3
  days to reach an area of 1000 cm$^2$.}
\end{enumerate}

\item \textbf{Conceptual.}

\ansline{For $g(x) = b^x$: as $x$ increases additively (by equal amounts), the output
multiplies by $b$ each time (multiplicative growth). For $f(x) = \log_b x$: as $x$
multiplies by $b$ (multiplicative input change), the output increases by 1 (additive growth).
The two patterns are exact ``reverses'' of each other because $f$ and $g$ are inverse
functions --- the roles of input and output are swapped.}

\item Let $h(x) = 2\log_3 x$.

\begin{enumerate}[leftmargin=1.5em, itemsep=6pt, label=\alph*.]
  \item $b = $ \ans{3} \qquad $a = $ \ans{2}

  \item $h(9) = 2\log_3 9 = 2 \cdot 2 = $ \ans{4} \qquad
        $h(27) = 2\log_3 27 = 2 \cdot 3 = $ \ans{6}

  \item \ansline{Every output of $h(x)$ is exactly twice the output of $\log_3 x$ for the
        same input. The vertical dilation by 2 stretches the outputs; the additive growth
        rate doubles (increases by 2 each time $x$ triples).}
\end{enumerate}

\item \textbf{AP-style MC.}

\begin{enumerate}[leftmargin=2em, itemsep=4pt, label=(\Alph*)]
  \item $g(x) = x^4$
  \item $g(x) = 4x$
  \item $g(x) = 4^x$ \quad \textcolor{keyred}{\textbf{$\leftarrow$ correct}}
  \item $g(x) = \log x$
\end{enumerate}

\begin{teachernote}
EK 2.10.A.4: $f(x) = \log_b x$ is the reflection of $g(x) = b^x$ over $y = x$, so the
inverse of $\log_4 x$ is $4^x$. (A) is a power function, not an exponential; (B) is linear;
(D) has a different base (base 10).
\end{teachernote}

\end{enumerate}

\vspace{0.12in}

\begin{reflectionbox}
What is one thing from today's lesson that you feel confident about, and one question you
still have about inverse functions or logarithms?

\writelines{2}
\end{reflectionbox}

\end{document}
